% --- Parte 3: Desarrollo de la Biblioteca ---
\section{Desarrollo de la Biblioteca}
\SectionFrame{Desarrollo de la Biblioteca}

\begin{frame}{OpalML}
    El objetivo de la libreria es abstraer la gestión del esquema CKKS en la \textbf{definición e inferencia de redes neuronales}. Ademas de aprovechar la capacidad de la GPU para procesar la inferencia.
    \vfill
    \textbf{Se abstrae del usuario}
    \begin{listaalineada}
        \setlength{\itemsep}{10pt}
        \item Las claves.
        \item Los niveles de los textos cifrados.
        \item Los objetos crudos del \textit{backend}.
        \item Gestión de niveles, ruido y \textit{bootstrapping}.
    \end{listaalineada}
    \vfill
\end{frame}

\begin{frame}{HEonGPU}
    \begin{columns}[t]
        \column{0.48\textwidth}
        \textbf{Por qué HEonGPU}
        \begin{itemize}
            \setlength{\itemsep}{8pt}
            \item La mayoria de librerias que implementan CKKS solo ejecutan sobre CPU.
            \item Entre las que usan el GPU, esta ofrece \textit{bootstrapping}.
            \item Esta en constante mantenimiento y con actualizaciones constantes.
            \item Es sencilla de integrar usando pybind11.
        \end{itemize}

        \column{0.48\textwidth}
        \textbf{Enlaces con pybind11}
        \begin{itemize}
            \setlength{\itemsep}{9pt}
            \item Contexto y parámetros del esquema.
            \item Clases para codificación, cifrado y descifrado.
            \item Clases para generación de claves.
            \item Objetos para realizar operaciones.
            \item Configuracion para el \textit{bootstrapping}.
        \end{itemize}
    \end{columns}
\end{frame}

\begin{frame}{Clases y funciones expuestas}
    El usuario solo interactúa con la interfaz pública; cada capa traduce las operaciones a las de la capa inferior hasta llegar a HEonGPU y a los núcleos CUDA.
    \vspace{0.15cm}
    \begin{center}
    \begin{tikzpicture}[
        scale=0.9, every node/.style={font=\scriptsize, transform shape},
        grupo/.style={draw=UCVBlue!45, rounded corners=3pt},
        caja/.style={draw=UCVBlue!60, rounded corners=2pt, fill=white, align=center,
                     inner sep=3pt},
        rotulo/.style={font=\tiny\bfseries, text=UCVBlue, anchor=north},
        flecha/.style={-{Stealth[length=5pt]}, very thick, UCVBlue!70},
    ]
        % --- Puntos de entrada ---
        \node[grupo, fill=UCVBlue!8, minimum width=3.4cm, minimum height=3.4cm] (ent) at (-4.85,0) {};
        \node[rotulo] at ([yshift=-4pt]ent.north) {PUNTOS DE ENTRADA (Python)};
        \node[caja, minimum width=2.9cm, minimum height=0.55cm] at (-4.85, 0.98) {Sequential};
        \node[caja, minimum width=2.9cm, minimum height=0.55cm] at (-4.85, 0.25) {Client};
        \node[caja, minimum width=2.9cm, minimum height=0.55cm] at (-4.85,-0.48) {FHEConfig};
        \node[caja, minimum width=2.9cm, minimum height=0.55cm] at (-4.85,-1.21) {FHEContext};

        % --- Componentes públicos ---
        \node[grupo, fill=UCVBlue!8, minimum width=5.3cm, minimum height=3.4cm] (comp) at (0.15,0) {};
        \node[rotulo] at ([yshift=-4pt]comp.north) {COMPONENTES PÚBLICOS};
        \node[caja, minimum width=4.9cm, minimum height=0.72cm] at (0.15, 0.90)
            {Capas\\ {\tiny Linear · Conv2D · ReLU · Square · Input}};
        \node[caja, minimum width=4.9cm, minimum height=0.72cm] at (0.15,-0.04)
            {Contenedores\\ {\tiny EncryptedVector · PlaintextVector · PlaintextTensor}};
        \node[caja, minimum width=4.9cm, minimum height=0.72cm] at (0.15,-0.98)
            {Configuración\\ {\tiny BootstrapConfig · SecurityLevel · config\_for\_circuit}};

        % --- Interno ---
        \node[grupo, fill=UCVBlue!16, minimum width=3.9cm, minimum height=3.4cm] (int) at (5.15,0) {};
        \node[rotulo] at ([yshift=-4pt]int.north) {INTERNO};
        \node[grupo, fill=UCVGrey!45, minimum width=3.5cm, minimum height=2.55cm] (enl) at (5.15,-0.28) {};
        \node[rotulo] at ([yshift=-4pt]enl.north) {CAPA DE ENLACES};
        \node[caja, minimum width=3.15cm, minimum height=0.6cm] at (5.15,-0.25)
            {Módulo de enlaces {\tiny (pybind11)}};
        \node[caja, minimum width=3.15cm, minimum height=0.6cm, fill=UCVGrey!75] at (5.15,-1.15)
            {HEonGPU {\tiny (C++/CUDA)}};

        \draw[flecha] (ent) -- (comp);
        \draw[flecha] (comp) -- (int);
    \end{tikzpicture}
    \end{center}
\end{frame}

\begin{frame}[fragile]{OpalML en uso}
\begin{lstlisting}[basicstyle=\ttfamily\fontsize{6.5}{7.1}\selectfont]
from opal_ml import Client, FHEContext, Sequential
from opal_ml.layers import Linear, ReLU
# Subconjunto representativo de los datos, del propio usuario.
from datos import datos_de_calibracion

# --- Servidor ---
model = Sequential([
    Linear(4, 8, W1, bias=b1),
    ReLU(),
    Linear(8, 2, W2, bias=b2),
])
config = model.generate_config(scale=40, relu_degrees=(3,))
ctx = FHEContext(config)                  # sin clave secreta
model.compile(ctx, datos_de_calibracion)  # calibra y prepara los pesos

# --- Cliente ---
client = Client(ctx)
ctx.set_client_params(client.key_params())   # claves de evaluacion
entrada_cifrada = client.encrypt(entrada)

# --- Servidor ---
salida_cifrada = model(entrada_cifrada)

# --- Cliente ---
resultado = client.decrypt(salida_cifrada)
\end{lstlisting}
\end{frame}

\begin{frame}{Multiplicación matriz-vector}
    Se empleó la multiplicación matriz-vector de \textbf{Halevi-Shoup} con una \textbf{descomposición} que reduce las operaciones.
    \vfill
    \begin{columns}[t]
        %% ================= 1. Método de las diagonales =================
        \column{0.46\textwidth}
        \textbf{1. Método de las diagonales}
        \vspace{0.45cm}

        \begin{tikzpicture}[
            scale=1.28, every node/.style={font=\tiny, transform shape},
            celda/.style={draw=UCVBlue!50, minimum width=0.30cm, minimum height=0.30cm,
                          inner sep=0pt, font=\tiny},
            xcelda/.style={celda, fill=white, minimum width=0.32cm},
            flecha/.style={-{Stealth[length=4pt]}, thick, UCVBlue!70},
        ]
            % --- Matriz W coloreada por diagonales ---
            \foreach \r in {0,1,2} {
                \foreach \c in {0,1,2} {
                    \pgfmathtruncatemacro{\d}{mod(\c-\r+3,3)}
                    \node[celda, fill={\ifcase\d UCVOrange!45\or UCVBlue!35\or UCVGrey!75\fi}]
                         at (\c*0.30, 0.77-\r*0.30) {};
                }
            }
            \node[font=\scriptsize] at (0.30, 1.12) {$W$};
            \draw[flecha] (0.9, 0.47) -- (1.32, 0.47);

            % --- Un producto por diagonal ---
            \foreach \i/\y in {0/0.90, 1/0.47, 2/0.04} {
                \node[anchor=east] at (1.77, \y) {$d_\i$};
                \foreach \c in {0,1,2} {
                    \node[celda, fill={\ifcase\i UCVOrange!45\or UCVBlue!35\or UCVGrey!75\fi}]
                         at (1.92+\c*0.30, \y) {};
                }
                \node at (2.83, \y) {$\odot$};
                \foreach \c in {0,1,2} {
                    \pgfmathtruncatemacro{\xi}{mod(\c+\i,3)}
                    \node[xcelda] at (3.15+\c*0.32, \y) {$x_\xi$};
                }
            }

            % --- La suma de los productos da el resultado ---
            \node[font=\scriptsize, anchor=east] at (1.8, -0.3) {$\sum$};
            \draw[UCVBlue!60, semithick] (1.77,-0.3) -- (3.95,-0.3);
            \node[anchor=east] at (1.77, -0.7) {$y$};
            \foreach \c in {0,1,2} {
                \node[celda, fill=UCVOrange!20] at (1.92+\c*0.30, -0.7) {};
            }
        \end{tikzpicture}

        %% ============ 2. Descomposición baby-step giant-step ============
        \column{0.54\textwidth}
        \textbf{2. Descomposición \textit{baby-step giant-step}}
        \vspace{0.5cm}

        \hspace*{0.5cm}
        \begin{tikzpicture}[
            scale=1.18, every node/.style={font=\tiny, transform shape},
            celda/.style={draw=UCVBlue!55, minimum width=0.80cm, minimum height=0.42cm,
                          inner sep=0pt, font=\tiny, fill=UCVBlue!10},
            blq/.style={draw=UCVBlue!60, rounded corners=2pt, fill=white,
                        minimum width=1.05cm, minimum height=0.42cm, inner sep=1pt, font=\tiny},
            flecha/.style={-{Stealth[length=4pt]}, thick, UCVBlue!70},
            giro/.style={-{Stealth[length=4pt]}, thick, UCVOrange},
            etq/.style={font=\tiny, text=UCVOrange, inner sep=1.5pt},
        ]
            % --- Matriz de diagonales: filas (bloques) por columnas ---
            \foreach \j in {0,1,2} {
                \foreach \k in {0,1,2} {
                    \pgfmathtruncatemacro{\idx}{3*\j+\k}
                    \node[celda] (C\j-\k) at (\k*0.88, 1.45-\j*0.62) {$d_{\idx}$};
                }
            }

            % --- Cada columna reutiliza una rotación más del vector ---
            \foreach \k [evaluate=\k as \n using int(\k+1)] in {0,1} {
                \draw[giro] (C0-\k.north) to[bend left=45] node[etq, above] {$+1$} (C0-\n.north);
            }

            % --- Cada fila da un bloque: solo productos y sumas ---
            \foreach \j in {0,1,2} {
                \node[blq] (B\j) at (3.30, 1.45-\j*0.62) {$\mathrm{bloque}_\j$};
                \draw[flecha, shorten >=2pt] (2.20, 1.45-\j*0.62) -- (B\j);
            }

            % --- Acumulación de abajo arriba: siempre la misma rotación ---
            \foreach \j [evaluate=\j as \p using int(\j-1)] in {2,1} {
                \draw[giro] (B\j.east) to[bend right=65] node[etq, right] {$+n_1$} (B\p.east);
            }

            \node[blq, fill=UCVOrange!25, minimum width=0.7cm] (Y) at (3.30, 2.16) {$y$};
            \draw[flecha] (B0) -- (Y);
        \end{tikzpicture}
    \end{columns}
    \vfill
\end{frame}

\begin{frame}{ReLU: aproximación y calibración}
    \begin{columns}[t]
        \column{0.50\textwidth}
        \textbf{Evaluación eficiente del polinomio}
        \vspace{0.1cm}
        \\ Cada componente es impar, de modo que se factoriza
        \[
            p(x) = x \cdot q(x^2), \qquad \deg q = n
        \]
        y se evalúa $q$ mediante el \textbf{esquema de Horner}.
        \begin{itemize}
            \item Profundidad: $n + 2$ niveles, frente a $2n+1$. \textbf{La mitad.}
        \end{itemize}

        \column{0.50\textwidth}
        \textbf{Calibración de la red}
        \vspace{0.1cm}
        \\ La aproximación solo es válida en $[-1,1]$.
        \begin{enumerate}
            \setlength{\itemsep}{5pt}
            \item Se propaga un conjunto representativo \textbf{en claro} y se registra el máximo absoluto $B$ de cada activación.
            \item Ese rango se \textbf{pliega en los pesos} de las capas vecinas.
        \end{enumerate}
    \end{columns}
    \vspace{0.5cm}
    \begin{center}
    \begin{tikzpicture}[every node/.style={font=\scriptsize},
        b/.style={draw=UCVBlue!60, rounded corners=2pt, minimum height=0.6cm, inner sep=4pt}]
        \node[b, fill=UCVBlue!10] (w1) at (-4.4,0) {$W_i / B$};
        \node[b, fill=UCVOrange!20] (r) at (-1.6,0) {ReLU en $[-1,1]$};
        \node[b, fill=UCVBlue!10] (w2) at (1.6,0) {$W_{i+1} \cdot B$};
        \draw[-{Stealth[length=4pt]}, thick, UCVBlue!70] (w1) -- (r);
        \draw[-{Stealth[length=4pt]}, thick, UCVBlue!70] (r) -- (w2);
        \node[anchor=west, align=left, font=\scriptsize] at (3.1,0) {$\mathrm{ReLU}(Bx) = B\,\mathrm{ReLU}(x)$\\ \textcolor{gray}{exacta, sin coste homomórfico}};
    \end{tikzpicture}
    \end{center}
\end{frame}

\begin{frame}{Automatización de los parámetros y validación}
    \begin{columns}[t]
        \column{0.50\textwidth}
        \textbf{Configuración derivada del circuito}
        \begin{itemize}
            \setlength{\itemsep}{6pt}
            \item El usuario indica el \textbf{factor de escala} y los \textbf{grados} de la aproximación de ReLU.
            \item La biblioteca deriva el resto de la configuración.
            \item Las \textbf{validaciones} rechazan una configuración inconsistente antes de llegar al \textit{backend}.
        \end{itemize}
        \vspace{0.3cm}
        \textbf{Desarrollo guiado por pruebas}
        \begin{itemize}
            \setlength{\itemsep}{6pt}
            \item Primero la interfaz, luego las pruebas, después la implementación.
            \item Cada operación cifrada reproduce su equivalente en claro.
        \end{itemize}

        \column{0.50\textwidth}
        \textbf{Gestión de niveles y \textit{bootstrapping}}
        \begin{itemize}
            \setlength{\itemsep}{6pt}
            \item Cada capa reporta su \textbf{profundidad multiplicativa}.
            \item En la compilación se suman y se comparan con los niveles disponibles.
            \item Si no alcanzan, se habilita el \textit{bootstrapping} en su variante \textbf{SLIM}.
            \item En inferencia el refresco se dispara \textbf{automáticamente}, incluso a mitad de una activación.
        \end{itemize}
    \end{columns}
\end{frame}
